% Örnek oyun — Dactile ile yazılmış, tamamen uydurma bir metin.
% Derleme: xelatex ornek-oyun.tex (iki kez; sahne kişi listeleri ikinci geçişte gelir)

\documentclass[draft]{tiyatro}

\title{Kar Masalı}
\author{Örnek Yazar}
\draft{0.1}
\copyrightnotice{2026 Örnek Yazar}
\address{İstanbul}
\email{ornek@example.com}

\karakter*[Selin][Köy öğretmeni]{SELİN}{se}
\karakter[Ozan Usta][Değirmenci, Selin'in amcası]{OZAN}{oz}
\karakter[Yaşlı bir yolcu]{İHTİYAR}{ih}
\karaktergrubu[Köylüler]{
	\gkarakter[Birinci Kadın]{KADIN}{ka}
	\gkarakter{ÇOBAN}{co}
}
\gizlikarakter{HERKES}{hep}
\setinlinename{oz}{Ozan Usta}

\begin{document}

\maketitle
\kisiler

\perde[Kar Yağarken]

\sahne[İç, Değirmen — Akşam]

\action{Işıklar yavaşça açılır. Değirmen taşı durmuştur.}

\stagedir{Karlı bir dağ köyü. Pencereden ışık sızar.}

\se Amca, yine mi taş durdu?
\oz Durdu kızım. \ic{elini taşa koyar} Soğuk her şeyi durdurur.
\se Ama \oz bu kış başka türlü diyordu.
\oz Ben öyle mi dedim? _Hiç_ hatırlamıyorum.

\ihdo Kapıyı çalar.

\se Kim o?
\ih Yolcuyum. Biraz ısınsam yeter.

\sedo\ozdo Birbirlerine bakarlar.

\middir{Rüzgârın sesi giderek artar.}

\sahne[Dış, Meydan — Gece]

\ka Kar dinmeyecek bu gece.
\co Sürü ağılda, merak etme.
\hep Hep birlikte: Dinmeyecek!
\se \ru{Добрый вечер}, dedi yolcu, \en{good evening}.
\beat

\perde[]

\sahne[][2.1]

\oz Sabah oldu.
\ih Ben gidiyorum artık.
\ihdo Çıkar.

\end{document}
